\section{Further reductions}\label{sec:other}

{This section presents the remaining algorithmic results of the introduction: directed unweighted APSP (Section~\ref{sec:dirapsp}), the real-valued versions of $3$SUM, APSP, and Exact Triangle (Section~\ref{sec:real}), the weighted $k$-Clique problems (Section~\ref{sec:kclique}), and the refutation of the rectangular hinted OMv conjectures of van den Brand, Nanongkai, and Saranurak~\cite{BNS19} (Section~\ref{sec:bns}). Each result but the last composes a known reduction with the algorithms of Sections~\ref{sec:reduction} and~\ref{sec:general}; the last applies the data structure of Section~\ref{sec:general} directly. We cite the reduction we use in each case and modify it only where our parameter regime requires it, saying what changes.}

\paragraph{What we use.}
We recall the three results of Sections~\ref{sec:reduction} and~\ref{sec:general} that the reductions plug into.
\begin{itemize}
\item \emph{Exact Triangle} (Section~\ref{sec:redproof}) asks whether a complete tripartite graph with parts $A,B,C$ of $n$ vertices and integer edge weights of absolute value at most $n^\nu$ has a triangle of weight zero. Theorem~\ref{thm:zwt} solves it deterministically in $O(n^{3-\varepsilon'}\log n)\leq O(n^{3-\varepsilon_T})$ time, where $\varepsilon'=0.00175$ and $\varepsilon_T=0.0017$, and Theorem~\ref{thm:3sumapsp} derives from this a deterministic $\tilde{O}(n^{3-\varepsilon'/3})$-time algorithm for the $(\min,+)$-product of two $n\times n$ matrices with integer entries of absolute value $n^{O(1)}$.
\item $\Lop(n,D)$ and $\#\Lop(n,D)$ (Definitions~\ref{def:lop} and~\ref{def:countlop}) are given a tripartite graph with parts $A$ and $B$ of $n$ vertices and a middle part of at most $D$ vertices, together with a set $W$ of \emph{query pairs} $(a,b)\in A\times B$, and ask for every query pair in $W$ whether $a$ and $b$ have a common neighbor in the middle part, or for the counting version, how many they have. Corollary~\ref{fact:lop} solves both deterministically in $O(|W|D^{0.437}+n^2/D^{0.063})$ time when $n\ge D^{18}$. Theorem~\ref{thm:red} reduces Exact Triangle to $\Lop$ deterministically by hashing the weights modulo a prime.
\item Corollary~\ref{cor:zt} is the data structure behind Corollary~\ref{fact:lop}. Given $X\in\Z^{N\times D}$ and $Y\in\Z^{D\times N}$ with $N\ge D^{18}$ and entries of absolute value $N^{O(1)}$, it preprocesses them deterministically in $O(N^2/D^{0.063})$ time, after which it computes any entry of $XY$ in $O(D^{0.437})$ time. Section~\ref{sec:real} applies it to matrices whose entries are not $0/1$, and Section~\ref{sec:bns} uses its queries online.
\end{itemize}
{The algorithms of Sections~\ref{sec:dirapsp}, \ref{sec:kclique}, and~\ref{sec:bns} are deterministic. Those of Section~\ref{sec:real} are randomized (Las Vegas), because the reductions of~\cite{CVX22} that they use as black boxes are randomized.}

\subsection{Directed APSP with small integer weights}\label{sec:dirapsp}

Let $\omega(1,r,1)$ denote the exponent of multiplying an $n\times n^{r}$ matrix by an $n^{r}\times n$ matrix, and let $\mu$ be the solution of $\omega(1,\mu,1)=1+2\mu$. Then $\mu\ge\frac12$, since $\omega(1,\mu,1)\ge2$, and $\mu<0.5275$ by the bounds of~\cite{ADVXXZ25}. Zwick's algorithm~\cite{Zwi02} solves APSP in directed unweighted graphs in $\tilde O(n^{2+\mu})$ time, and the Directed Unweighted APSP hypothesis~\cite{CVX21,CVX23,Fis26} asserts that no algorithm is polynomially faster. We refute this in Theorem~\ref{thm:dirapsp}.

We revisit Theorem~3.1 of Chan, Vassilevska Williams, and Xu~\cite{CVX21}, which is just a presentation of Zwick's algorithm~\cite{Zwi02} as a reduction. Zwick's algorithm is typically used in its randomized form because the proof is simple, so the reduction in \cite{CVX21} was stated in its randomized form. However, following an argument by Yuster and Zwick~\cite[Section~8]{YZ05b}, the reduction has a deterministic version as well:

\begin{theorem}[Zwick's algorithm as a deterministic reduction~\cite{Zwi02,CVX21}]\label{thm:detred}
If the $(\min,+)$-product of an $n\times n^{\mu}$ matrix by an $n^{\mu}\times n$ matrix, both with integer entries of absolute value at most $\tilde{O}(n^{1-\mu})$, can be computed in $O(n^{2+\mu-\varepsilon_1})$ time for a constant $\varepsilon_1>0$, then APSP in directed $n$-vertex graphs with integer weights of absolute value $\polylog n$ and no negative cycles can be solved in $O(n^{2+\mu-\varepsilon''})$ time for a constant $\varepsilon''>0$ that depends on $\varepsilon_1$ and $\mu$. If the algorithm for the product is deterministic, then so is the algorithm for APSP.
\end{theorem}


The randomized proof of~\cite{CVX21} follows the randomized version of Zwick's algorithm, which takes random samples of vertices to hit the long shortest paths, so the algorithm that it produces is randomized. {Zwick~\cite[Sections~6 and~7]{Zwi02} also gives a deterministic version of his algorithm, which constructs the hitting sets (his bridging sets) from witnesses of the products by a greedy hitting-set computation, within the same running time up to polylogarithmic factors, and Yuster and Zwick~\cite[Section~8]{YZ05b} announce, without details, a faster deterministic construction of bridging sets for their distance oracle. Zwick's construction lists a path with $s$ vertices for each of the $n^2$ pairs, which is more than the reduction can afford when $s$ is large. We therefore give a deterministic version of the reduction, which lists only walks with an endpoint in the previous hitting set, in the spirit of~\cite{YZ05b}. We believe that this is likely what Yuster and Zwick~\cite{YZ05b} had in mind.} 


\begin{proof}[Proof of Theorem~\ref{thm:detred}]
The deterministic algorithm follows the randomized reduction, with two changes. It constructs the hitting sets deterministically, by a hierarchy of bridging sets in the manner of Zwick~\cite[Section~6]{Zwi02}, and it computes the distances that this construction needs only to and from small sets of vertices, like the distance oracle of Yuster and Zwick~\cite{YZ05b}. The hypothesized algorithm for the $(\min,+)$-product is used only for the values of its products. (This is convenient but not essential: the witness-finding method of Seidel~\cite{Sei95} and Alon, Galil, Margalit, and Naor~\cite{AGMN92,GM93}, derandomized by Alon and Naor~\cite{AN96} and applied to $(\min,+)$-products by Zwick~\cite[Lemma~3.3]{Zwi02}, uses the product algorithm only on submatrices, so any deterministic algorithm for the product can be made to return witnesses at a polylogarithmic cost.) The price is a factor $n^{\gamma}$ in the running time, for an arbitrarily small constant $\gamma>0$, which the polynomial saving of the hypothesis absorbs. Throughout, $M$ is the largest absolute weight, which is polylogarithmic, and $\tilde O$ hides polylogarithmic factors.

\emph{Bridges.} For vertices $u$ and $v$, let $\delta(u,v)$ be the distance from $u$ to $v$, and let $\eta(u,v)$ be the smallest number of edges on a shortest path from $u$ to $v$. Every subpath of a shortest path with the fewest edges is again a shortest path with the fewest edges. For a threshold $t$ and a constant $C$, a set $B\subseteq V$ is a \emph{$(t,C)$-bridge} if every pair $(u,v)$ with $\eta(u,v)\ge t$ has a shortest \emph{walk} from $u$ to $v$ that passes through a vertex of $B$ and has at most $C\eta(u,v)$ edges. This is Zwick's strong bridging set with a constant-factor slack in the number of edges, and with walks in place of paths because a shortest walk may run around zero-weight cycles. Let $R:=\lceil n^{\gamma}\rceil$ (so $R$ is tiny!) and $t_i:=\min\{R^{i},n\}$ for $i=0,1,\dots,\lceil1/\gamma\rceil$. We shall construct $(t_i,C_i)$-bridges $B_i$ with $|B_i|=\tilde O(n/t_i)$, starting with $B_0=V$ and $C_0=1$, where $C_i$ is a constant that depends only on $i$. The slack $C_i$ grows by a constant factor from each level to the next. This is why the levels are a factor $n^{\gamma}$ apart, so that there are only $O(1/\gamma)$ of them, rather than a factor $3/2$ apart as in Zwick's algorithm.

\emph{The tables $E_i(T,H)$.} The distances to and from small sets of vertices are computed by the following routine. Given the bridges $B_0,\dots,B_i$, a set $T$ of $\tilde O(n/t_i)$ \emph{target} vertices, and a \emph{horizon} $H$ with $t_i\le H\le O(Rt_i)$, the routine returns two tables, $E_i(T,H)[V,T]$ and $E_i(T,H)[T,V]$, with three properties: every finite entry is the weight of a walk that the routine represents explicitly, as described next; every represented walk has $O(H)$ edges, with a constant that depends on $i$; and the entry for $(u,v)$ is $\delta(u,v)$ whenever $\eta(u,v)\le H$. The products in the routine are computed with the deterministic $(\min,+)$-product algorithm with witnesses of Zwick~\cite[Lemma~3.3]{Zwi02}, which finds the witnesses deterministically by the method of Alon and Naor~\cite{AN96}. A witness of an entry $(u,v)$ of a product is an index $k$ at which its minimum is attained, and we store it with the entry. The walk that the entry represents is then the concatenation of the walks represented by the entries $(u,k)$ and $(k,v)$ of the two factors, and it is listed, when needed, by following the stored witnesses recursively down to the entries of the weight matrix, which represent single edges; this takes $O(H)$ time, because the walk is a concatenation of $O(H)$ such entries. For a product in which one dimension is $n$ and the other two are at most $q$, with entries of absolute value at most $W$, that algorithm takes $\tilde O(Wnq^{\omega-1})$ time, by cutting the product into $n/q$ square products.

The routine for level $0$ takes the weight matrix with zeroes on the diagonal and squares it $O(\log H)$ times, with the entries larger than $2HM$ replaced by $\infty$, and keeps the rows and columns of $T$. The routine for level $i\ge1$ first calls the routine for level $i-1$ with the target set $T\cup B_i$ and the horizon $L:=2C_it_i$. This call is okay, because $|T\cup B_i|=\tilde O(n/t_i)\le\tilde O(n/t_{i-1})$ and $L=O(Rt_{i-1})$. Let $F$ be the pair of tables that it returns; they have exact distances for the pairs with $\eta\le L$, and they have the entries of $V\times(T\cup B_i)$ and $(T\cup B_i)\times V$. Let $Q$ be the smallest power of two above $H/t_i$, and set the diagonal of $F[B_i,B_i]$ to zero. The routine computes, by repeated squaring,
\[
E_i(T,H)[V,T]:=\min\bigl\{F[V,T],\;F[V,B_i]\star F[B_i,B_i]^{\star Q}\star F[B_i,T]\bigr\},
\]
where $\star$ is the $(\min,+)$-product and the minimum is entrywise, and it computes $E_i(T,H)[T,V]$ symmetrically.

The entry for a pair $(u,v)$ with $v\in T$ and $\eta(u,v)\le H$ is exact. To see this, take a shortest path from $u$ to $v$ with the fewest edges, and cut it into blocks of $t_i$ edges followed by a remainder of fewer than $t_i$ edges. Each block is a shortest path with the fewest edges between its endpoints, so the bridge property of $B_i$ replaces it by a shortest walk between the same endpoints that has at most $C_it_i$ edges and passes through a vertex of $B_i$; choose one such vertex in each block. The result is a shortest walk from $u$ to $v$ through at most $H/t_i$ chosen vertices of $B_i$, in which $u$, the chosen vertices, and $v$ are consecutively at most $L=2C_it_i$ edges apart. The table $F$ is exact on these consecutive pairs, so the product above finds $\delta(u,v)$. The represented walks have $O(H)$ edges, since those of $F$ have $O(L)$ edges and at most $Q+2=O(H/t_i)$ of them are concatenated. Every product of the routine has one dimension $n$, the two others $\tilde O(n/t_i)$, and entries $O(HM)=\tilde O(Rt_i)$, so it takes $\tilde O(Rt_i\cdot n\cdot(n/t_i)^{\omega-1})\le\tilde O(n^{\omega+\gamma})$ time, and the recursion has $O(1/\gamma)$ levels.

\emph{The stages.} The algorithm maintains a matrix of distance estimates, initially the weight matrix with zeros on the diagonal, and runs one stage for each level $i=0,1,\dots$ until $t_{i+1}=n$. The stage for level $i$ computes the tables $E_i:=E_i(B_i,C_it_{i+1})$, then the $(\min,+)$-product
\[
E_i[V,B_i]\star E_i[B_i,V],
\]
and takes its entrywise minimum with the matrix of estimates. Then, if $t_{i+1}<n$, it constructs $B_{i+1}$ from the walks represented in $E_i$, as described below. The product fixes every pair $(u,v)$ with $t_i\le\eta(u,v)\le t_{i+1}$: the bridge property of $B_i$ gives $b\in B_i$ with $\delta(u,v)=\delta(u,b)+\delta(b,v)$ and $\eta(u,b),\eta(b,v)\le C_i\eta(u,v)\le C_it_{i+1}$, so both $E_i[u,b]$ and $E_i[b,v]$ are exact. The stages thus find all distances, and every estimate is the weight of a walk, so none is too small.

\emph{The cost of the stages.} Let $s:=t_{i+1}$. The product of the stage for level $i$ has middle dimension $|B_i|=\tilde O(n/t_i)=\tilde O(Rn/s)$ and entries $\tilde O(s)$. Cutting $B_i$ into $\tilde O(R)$ groups of $\lceil n/s\rceil$ vertices turns it into $\tilde O(n^{\gamma})$ products of an $n\times\lceil n/s\rceil$ matrix by an $\lceil n/s\rceil\times n$ matrix with entries $\tilde O(s)$. Up to this grouping, these are the products of the randomized reduction, and the analysis of~\cite[Theorem~3.1]{CVX21} applies to them. Fix a small constant $a>0$. When $s\le n^{1-\mu-a}$, the middle dimension $n^{r}:=\lceil n/s\rceil$ has $r\ge\mu+a$, and fast rectangular matrix multiplication computes the product in $\tilde O(s\,n^{\omega(1,r,1)})=O(n^{2+\mu-ca})$ time~\cite[Lemma~3.3]{Zwi02}, for a constant $c>0$, by the convexity of $\omega(1,\cdot,1)$, since $\omega(1,\mu,1)=1+2\mu$ and $\omega<2+\mu$. When $s\ge n^{1-\mu+a}$, the naive algorithm computes it in $O(n^{2}\cdot n/s)=O(n^{2+\mu-a})$ time. In between, the middle dimension is at most $n^{\mu+a}$ and the entries are $\tilde O(n^{1-\mu+a})$; we cut the middle dimension into at most $n^{a}$ pieces of $n^{\mu}$ columns and pad each piece to an instance with $n^{1+O(a)}$ rows, so that its entries fit the range of the hypothesis, and the hypothesized algorithm computes it in $O(n^{2+\mu-\varepsilon_1+O(a)})$ time. With $a$ small enough, the product of every stage takes $O(n^{2+\mu+\gamma-\varepsilon_2})$ time, for a constant $\varepsilon_2>0$ that depends on $\varepsilon_1$ and $\mu$.

\emph{The next bridge.} Let $s:=t_{i+1}<n$. The bridge $B_{i+1}$ is read off the walks represented in $E_i=E_i(B_i,C_is)$, and this is where the small target set pays off. These are $2n|B_i|$ walks with $O(s)$ edges each, so listing them all takes $\tilde O(n|B_i|s)=\tilde O(n^{2}R)\leq \tilde O(n^{2+\gamma})$ time, whereas listing a path for every pair of vertices, as Zwick's construction does, would take $O(n^{2}s)$ time, which exceeds our budget at the large scales. From each listed walk, erase the cycles, and keep the resulting path if it has at least $s/2$ edges. Let $B_{i+1}$ be a greedy hitting set of the vertex sets of the kept paths, each of which has more than $s/2$ vertices. It has $O((n/s)\log n)$ vertices by the analysis of the greedy heuristic~\cite{Lov75,Chv79}, and it is computed in time linear in the total size of the sets, up to a logarithmic factor.

This $B_{i+1}$ is an $(s,C_{i+1})$-bridge for a constant $C_{i+1}$ that depends only on $i$. Consider first a pair $(x,y)$ with $\eta(x,y)=s$. The bridge $B_i$ gives $b$ with $\delta(x,y)=\delta(x,b)+\delta(b,y)$ and $\eta(x,b)+\eta(b,y)\le C_is$, so the entries of $E_i$ for $(x,b)$ and $(b,y)$ are exact, and the two walks that represent them are shortest walks. A cycle on a shortest walk has weight zero, so the two paths obtained by erasing the cycles are still shortest, and their concatenation is a shortest walk from $x$ to $y$ with $O(s)$ edges. This walk has at least $s$ edges, since otherwise erasing its cycles would give a shortest path from $x$ to $y$ with fewer than $\eta(x,y)$ edges. So one of the two paths has at least $s/2$ edges and was kept, and $B_{i+1}$ contains one of its vertices, which lies on the concatenated walk. For a pair with $\eta(x,y)>s$, apply this to the first $s$ edges of a shortest path from $x$ to $y$ with the fewest edges, and keep the rest of that path; the walk obtained has at most $C_{i+1}s+\eta(x,y)-s\le C_{i+1}\eta(x,y)$ edges.

\emph{Running time.} There are $O(1/\gamma)$ stages. Each computes its tables in $\tilde O(n^{\omega+\gamma})$ time, its product in $O(n^{2+\mu+\gamma-\varepsilon_2})$ time, and its bridge in $\tilde O(n^{2+\gamma})$ time. Choosing $\gamma$ smaller than both $\varepsilon_2/2$ and $(2+\mu-\omega)/2$ gives the running time $O(n^{2+\mu-\varepsilon''})$ for a constant $\varepsilon''>0$, and every step is deterministic.
\end{proof}



We compose the known reduction with the deterministic $(\min,+)$-product algorithm of Theorem~\ref{thm:3sumapsp}.

\begin{theorem}[Directed unweighted APSP]\label{thm:dirapsp}
There is a constant $\varepsilon''>0$ such that APSP in directed unweighted $n$-vertex graphs, and more generally in directed $n$-vertex graphs with integer weights of absolute value at most $(\log n)^c$ and no negative cycles, for any constant $c$, can be solved in $O(n^{2+\mu-\varepsilon''})$ time by a deterministic algorithm.
\end{theorem}
\begin{proof}
Cut the $n\times n^{\mu}$ matrix into $n^{1-\mu}$ blocks of $n^{\mu}$ consecutive rows, and the $n^{\mu}\times n$ matrix into $n^{1-\mu}$ blocks of $n^{\mu}$ consecutive columns. The $(\min,+)$-product then consists of $n^{2-2\mu}$ products of two $n^{\mu}\times n^{\mu}$ matrices, one for each pair of blocks. Their entries have absolute value $\tilde{O}(n^{1-\mu})\leq\tilde{O}(n^{\mu})$, because $\mu\ge\frac12$, so Theorem~\ref{thm:3sumapsp} computes each of them in $\tilde{O}((n^{\mu})^{3-\varepsilon'/3})$ time, and all of them in $\tilde{O}(n^{2+\mu-\mu\varepsilon'/3})$ time. This is $O(n^{2+\mu-\varepsilon_1})$ for any constant $\varepsilon_1<\mu\varepsilon'/3$, so~\cite[Theorem~3.1]{CVX21} applies, in the deterministic form of Theorem~\ref{thm:detred}, since the algorithm of Theorem~\ref{thm:3sumapsp} is deterministic.
\end{proof}

\subsection{3SUM, APSP, and Exact Triangle with real inputs}\label{sec:real}

\begin{theorem}[Real inputs]\label{thm:real}
In the real RAM in which the only operations allowed on real numbers are comparisons, additions, and subtractions, Las Vegas algorithms solve the following problems on real inputs: $3$SUM on $n$ numbers in $O(n^{1.998})$ expected time, and the $(\min,+)$-product of two $n\times n$ matrices, APSP with no negative cycles, and Exact Triangle in $O(n^{2.998})$ expected time. The bounds also hold with probability $1-n^{-c}$ for any constant $c$.
\end{theorem}

The reductions are by Chan, Vassilevska Williams, and Xu~\cite{CVX22}, who reduce the real-valued problems to the counting version of All-Edges Sparse Triangle on graphs with $m$ edges (and real APSP also to the Boolean version). Our version of All-Edges Sparse Triangle in Corollary~\ref{fact:lop} needs lopsided sparse graphs.
To deal with this slight change, we follow the reductions of~\cite{CVX22} down to the point where the All-Edges Sparse Triangle oracle is called, and change the call. At that point, the reductions only need the oracle to compute \emph{comparison counts}, which we define next. These come from Fredman's trick~\cite{Fre76}: $a'+b'<a+b$ if and only if $a'-a<b-b'$. This turns a comparison of two sums into a comparison of a number that depends only on the row against a number that depends only on the column.

{We first define comparison counts, then restate the reductions of~\cite{CVX22} as reductions to comparison counts (Lemma~\ref{lem:realred}), and then solve the comparison counts problem: following~\cite{CVX22}, Lemma~\ref{lem:cmp} reduces it to the wanted entries of one thin matrix product using Matou\v sek's technique for the dominance product~\cite{Mat91}, and Corollary~\ref{cor:cmp} applies Corollary~\ref{cor:zt} to that product.} The proof of Theorem~\ref{thm:real} then adds up the running times. The only randomized steps are in the reductions of~\cite{CVX22}, and the algorithms never output a wrong answer, so they are Las Vegas.

\paragraph{Comparison counts.}
We are given $n$ \emph{row lists}: for every $r\in[n]$, a list $\mathcal R_r$ of at most $d$ real numbers $x$, where each $x\in\mathcal R_r$ has a \emph{color} $\chi(x)\in[d]$ associated with it.

Similarly, we are given $n$ \emph{column lists}: for every $c\in[n]$, a list $\mathcal C_c$ of at most $d$ real numbers $y$, each $y\in\mathcal C_c$ again with a color $\chi(y)\in[d]$ associated with it.

Given a set $P\subseteq[n]^2$ of (row, column) pairs, we want, for every $(r,c)\in P$, the number
\[
\gamma(r,c):=\big|\{(x,y):\ x\in\mathcal R_r,\ y\in\mathcal C_c,\ \chi(x)=\chi(y),\ x<y\}\big| .
\]

\paragraph{The reductions.}
We follow the reductions of~\cite{CVX22} down to their two counting problems~\cite[Problems~4.1 and~4.5]{CVX22}, which we define in the proof of Lemma~\ref{lem:realred} below, and turn every call of a counting problem into comparison counts, in place of the triangle-counting instances of~\cite[Lemmas~4.3 and~4.7]{CVX22}. Throughout, $d\in[n]$ is a parameter, which the proof of Theorem~\ref{thm:real} chooses.

\begin{lemma}[after Chan, Vassilevska Williams, and Xu~\cite{CVX22}]\label{lem:realred}
Let $1\le d\le n$.
\begin{enumerate}
\item[(a)] The $(\min,+)$-product of two $n\times n$ real matrices, APSP with real edge weights and no negative cycles on $n$ vertices, and Exact Triangle with real weights on $n$ vertices reduce, with Las Vegas randomization, to $\tilde{O}(n/d)$ \emph{counting calls} and $\tilde{O}(n^3/d)$ further time. A counting call consists of $d$ comparison counts, each with $n$ row lists and $n$ column lists of at most $d$ numbers, with at most one number of each color in a list, and with pair sets $P_1,\dots,P_d$ satisfying $\sum_k|P_k|=n^2$; forming its lists costs $O(nd^2)$ subtractions.
\item[(b)] $3$SUM on $n$ real numbers reduces, with Las Vegas randomization, to polylogarithmically many comparison counts, each with $n$ row lists and $n$ column lists of at most $d$ numbers, all of the same color, and a pair set $P$ of size $O(n^2/d)$, and $\tilde{O}(n^2/d)$ further time; forming the lists of a comparison count costs $O(nd)$ subtractions.
\end{enumerate}
A comparison count may have $\le$ in place of $<$. The only operations on real numbers performed by the reductions are comparisons, additions, and subtractions.
\end{lemma}

\begin{proof}
\emph{(a) $(\min,+)$-product, APSP, and Exact Triangle.} Let $A,B,C\in\mathbb R^{n\times n}$. As in~\cite[Section~3.2]{CVX22}, cut $A$ into $n/d$ blocks $A'\in\mathbb R^{n\times d}$ of $d$ consecutive columns, and $B$ into the corresponding blocks $B'\in\mathbb R^{d\times n}$ of $d$ consecutive rows. For each pair $(A',B')$ we solve~\cite[Problem~3.2]{CVX22}: find, for every $(i,j)$, the predecessor and the successor of $C[i,j]$ among the $d$ sums $A'[i,k]+B'[k,j]$, where a sum equal to $C[i,j]$ counts as its predecessor. For Exact Triangle, $A[i,k]:=w(i,k)$, $B[k,j]:=w(k,j)$, and $C[i,j]:=-w(i,j)$, and there is a triangle of weight zero if and only if, in some block, some $C[i,j]$ is its own predecessor. For the $(\min,+)$-product of $A$ and $B$, we take every $C[i,j]$ smaller than all the sums, for instance $C[i,j]:=\min_k A[i,k]+\min_k B[k,j]-1$. Then the successor of $C[i,j]$ in a block is the smallest sum of the block, and the product is the entrywise minimum over the $n/d$ blocks, which costs $O(n^3/d)$ comparisons. APSP with real weights and no negative cycles can be computed by successive squaring of the weight matrix, performing $\lceil\log_2n\rceil$ such products.

Fix a pair $(A',B')$.  We solve the following \emph{counting problem}~\cite[Problem~4.1]{CVX22}: we are given a \emph{pivot} $k_{ij}\in[d]$ for every $(i,j)\in[n]^2$ and a subset $S\subseteq[d]$, and we must count, for every $(i,j)$, the indices $k'\in S$ with $A'[i,k']+B'[k',j]<A'[i,k_{ij}]+B'[k_{ij},j]$; a call of the problem may also have $\le$ in place of $<$. 

By Lemmas~3.4 and~4.2 of~\cite{CVX22}, the problem for one pair $(A',B')$ reduces, with Las Vegas randomization, to polylogarithmically many calls of the counting problem, and $\tilde{O}(n^2)$ further time: Lemma~3.4 finds the predecessors and successors by a randomized search that repeatedly asks, for every $(i,j)$, for a sum strictly between two given sums, and Lemma~4.2 finds these sums by counting, for each of the two given sums, the sums below it, and by sampling at random. Lemma~4.3 of~\cite{CVX22} turns a call into $O(d)$ triangle-counting instances; we turn it into $d$ comparison counts instead. 

Following~\cite{CVX22}, we group the pairs by their pivot, $P_k:=\{(i,j):k_{ij}=k\}$. For $k\in[d]$, let the row list $\mathcal R_i$ consist of the numbers $A'[i,k']-A'[i,k]$, $k'\in S$, each with color $k'$, and let the column list $\mathcal C_j$ consist of the numbers $B'[k,j]-B'[k',j]$, $k'\in S$, again with color $k'$.

By Fredman's trick, $A'[i,k']+B'[k',j]<A'[i,k]+B'[k,j]$ if and only if $A'[i,k']-A'[i,k]<B'[k,j]-B'[k',j]$, where both numbers in the latter comparison have color $k'$.

Thus the count of a pair $(i,j)\in P_k$ is the comparison count $\gamma(i,j)$ of these lists with $P:=P_k$. Forming the lists costs us $O(nd)$ subtractions for each $k$. Over the $n/d$ pairs $(A',B')$, and the $\lceil\log_2n\rceil$ products in the case of APSP, this gives $\tilde{O}(n/d)$ counting calls and $\tilde{O}(n^3/d)$ additional time, which includes the $O(n^3/d)$ comparisons of the entrywise minimum.

\emph{(b) $3$SUM.} Let $A,B,C$ be sets of $n$ reals, the problem asks whether $a+b=c$ for some $a\in A$, $b\in B$, and $c\in C$. The $3$SUM version with one set and $a+b+c=0$ reduces to this version in the standard way.

Now, sort $A$ and $B$, and cut them into $n/d$ consecutive blocks $A_1,\dots,A_{n/d}$ and $B_1,\dots,B_{n/d}$ of $d$ numbers each, writing $A_i[k]$ for the $k$-th number of $A_i$. Here the counting problem~\cite[Problem~4.5]{CVX22} is the following: we are given a set $\mathcal Q$ of $O(n^2/d)$ quadruples $(i,j,k,\ell)$, and we must count, for every $(i,j,k,\ell)\in\mathcal Q$, the pairs $(k',\ell')\in[d]^2$ with $A_i[k']+B_j[\ell']<A_i[k]+B_j[\ell]$; again, a call may have $\le$ in place of $<$, and it may restrict $k'$ and $\ell'$ to subsets of $[d]$. By Section~3.3 and Lemmas~3.9 and~4.6 of~\cite{CVX22}, $3$SUM reduces, with Las Vegas randomization, to polylogarithmically many calls of this counting problem, and $\tilde{O}(n^2/d)$ additional time. By Fredman's trick, the condition is $A_i[k']-A_i[k]<B_j[\ell]-B_j[\ell']$. So a call is one comparison count: the rows are the $n$ pairs $(i,k)$ and the columns are the $n$ pairs $(j,\ell)$, the row list $\mathcal R_{(i,k)}$ consists of the $d$ numbers $A_i[k']-A_i[k]$, $k'\in[d]$, the column list $\mathcal C_{(j,\ell)}$ of the $d$ numbers $B_j[\ell]-B_j[\ell']$, $\ell'\in[d]$, all numbers have the same color, and $P:=\mathcal Q$. A restriction of $k'$ or of $\ell'$ deletes numbers from the lists. Forming the lists can be done using $O(nd)$ subtractions.
\end{proof}

\paragraph{Solving comparison counts.}
In the comparison counts of Lemma~\ref{lem:realred}(a), every list has at most one number of each color, so that $\gamma(r,c)$ is the number of colors $k$ with $A[r,k]<B[k,c]$, where $A[r,k]$ is the number of color $k$ in $\mathcal R_r$ and $B[k,c]$ the one in $\mathcal C_c$ (a color missing from either list is skipped); this is the dominance product of the $n\times d$ matrix $A$ and the $d\times n$ matrix $B$ at the positions in $P$. In those of Lemma~\ref{lem:realred}(b), all numbers have the same color, and $\gamma(r,c)$ is the number of pairs $x<y$ between the two lists. Merging the two sorted lists of every pair in $P$ takes $O(|P|d)$ time. 
To improve on this, we utilize Matou\v sek's technique for the dominance product~\cite{Mat91}, which counts the pairs that lie in different blocks of the sorted order by one thin matrix product, and the pairs in the same block directly.

\begin{lemma}[after Matou\v sek~\cite{Mat91}]\label{lem:cmp}
Let $s\ge1$ be an integer and $D'':=\lceil2nd/s\rceil+d$. Given row lists $\mathcal R_r$ and column lists $\mathcal C_c$ as above and a set $P\subseteq[n]^2$, we can build matrices $X\in\{0,\dots,d\}^{n\times D''}$ and $Y\in\{0,\dots,d\}^{D''\times n}$ and compute integers $\gamma_2(r,c)$, $(r,c)\in P$, with
\[
\gamma(r,c)=(XY)[r,c]+\gamma_2(r,c)\qquad\text{for every }(r,c)\in P,
\]
deterministically in $O(nD''+(|P|+nds+ds^2)\log n)$ time. The only operations on real numbers are the comparisons made when sorting the numbers of each color. The same holds with $\le$ in place of $<$ in the definition of $\gamma$.
\end{lemma}

\begin{proof}
For each color $k$, sort the numbers of color $k$ of all the lists, and break ties so that, among equal numbers, those from column lists precede those from row lists. Let $\mathrm{pos}(x)$ be the position of $x$ in this order. For $x$ from a row list and $y$ from a column list of the same color, $x<y$ holds if and only if $\mathrm{pos}(x)<\mathrm{pos}(y)$. (For $\le$, we reverse the tie-breaking rule.) Cut the order into \emph{blocks} of $s$ consecutive positions, and let $\mathrm{blk}(x):=\lfloor\mathrm{pos}(x)/s\rfloor$. Then $x<y$ holds if and only if either $\mathrm{blk}(x)<\mathrm{blk}(y)$, or $\mathrm{blk}(x)=\mathrm{blk}(y)$ and $\mathrm{pos}(x)<\mathrm{pos}(y)$. As in Matou\v sek's algorithm, we count the pairs of the first kind, which lie in different blocks, by a matrix product, and the pairs of the second kind, which lie in the same block, by brute force.

\emph{Different blocks.} The lists contain at most $2nd$ numbers in all, so there are at most $2nd/s+d\le D''$ pairs (color, block), and we index the columns of $X$ and the rows of $Y$ by the pairs, padding with zeros if needed. Let $X[r,(k,\beta)]$ be the number of numbers of color $k$ in $\mathcal R_r$ that lie in block $\beta$, and let $Y[(k,\beta),c]$ be the number of numbers of color $k$ in $\mathcal C_c$ that lie in a block after $\beta$. Then $(XY)[r,c]=\sum_{(k,\beta)}X[r,(k,\beta)]\,Y[(k,\beta),c]$ is the number of pairs $(x,y)$ of the same color with $\mathrm{blk}(x)<\mathrm{blk}(y)$, as required, and the entries of $X$ and $Y$ are between $0$ and $d$ because a list has at most $d$ numbers. We use $O(nd\log n)$ comparisons to sort and then fill in $X$ and $Y$ in 
 $O(nD'')$ time.

\emph{Same block.} For every color $k$ and block $\beta$, let $I$ be the numbers $x$ of row lists in the block, each with its row, and let $J$ be the numbers $y$ of column lists in the block, each with its column, so that $|I|+|J|\le s$. For every $(x,y)\in I\times J$ with $\mathrm{pos}(x)<\mathrm{pos}(y)$ whose row and column form a pair $(r,c)\in P$ (this can be tested via a binary search in $P$ after it is sorted), we add $1$ to $\gamma_2(r,c)$. Then $\gamma_2$ counts the pairs of the second kind. Since $|I||J|\le s^2/4$ and there are at most $2nd/s+d$ blocks, we enumerate at most $(2nd/s+d)s^2/4\le nds+ds^2$ pairs $(x,y)$. Together with sorting $P$ and outputting the $|P|$ integers $\gamma_2(r,c)$, this takes $O((|P|+nds+ds^2)\log n)$ time.
\end{proof}

\begin{corollary}\label{cor:cmp}
Let $d\le n^{1/40}$, and let $n$ be larger than a suitable constant. Given row lists $\mathcal R_r$ and column lists $\mathcal C_c$ as above and a set $P\subseteq[n]^2$, the counts $\gamma(r,c)$, $(r,c)\in P$, with $<$ or with $\le$, can be computed deterministically in
\[
O\big(|P|\,n^{0.0229}+n^2/d^{0.126}+n^{2-1/440}\log n\big)
\]
time, by an algorithm whose only operations on real numbers are comparisons.
\end{corollary}
\begin{proof}
Apply Lemma~\ref{lem:cmp} with $s:=\lceil n^{1-1/440}/d\rceil$, so that $D''\le2d^2n^{1/440}+d+1\le3d^2n^{1/440}$, and pad the product to the inner dimension $D^*:=\lceil3d^2n^{1/440}\rceil$ with zero columns of $X$ and zero rows of $Y$. 
Since $d\le n^{1/40}$, we have $D^*\le4n^{23/440}$, so $(D^*)^{18}\le4^{18}n^{0.941}\le n$, and Corollary~\ref{cor:zt} with $N:=n$ and $D:=D^*$ computes $(XY)[r,c]$ for all $(r,c)\in P$ deterministically in $O(|P|(D^*)^{0.437}+n^2/(D^*)^{0.063})$ time. Here $(D^*)^{0.437}\le4^{0.437}n^{0.437\cdot23/440}=O(n^{0.0229})$ and $(D^*)^{0.063}\ge(3d^2)^{0.063}\ge d^{0.126}$. 

Since $nD^*=O(n^{1.06})$, $nds\le2n^{2-1/440}$ and $ds^2\le4n^{2-1/440}$, the time of Lemma~\ref{lem:cmp} itself is $$O(nD^*+(|P|+nds+ds^2)\log n)=O(|P|\,n^{0.0229}+n^{2-1/440}\log n).$$
\end{proof}

\begin{proof}[Proof of Theorem~\ref{thm:real}]
Let $d:=\lfloor n^{1/40}\rfloor$, apply Lemma~\ref{lem:realred} with this $d$, and answer every comparison count by Corollary~\ref{cor:cmp}.

\emph{$(\min,+)$-product, APSP, and Exact Triangle.} By Corollary~\ref{cor:cmp}, the $d$ comparison counts of a counting call cost
\[
\sum_{k\in[d]}O\big(|P_k|\,n^{0.0229}+n^2/d^{0.126}+n^{2-1/440}\log n\big)=O(n^{2.0229}\log n),
\]
since $\sum_k|P_k|=n^2$, $d\cdot n^2/d^{0.126}\le n^{2+0.874/40}\le n^{2.0219}$, and $d\cdot n^{2-1/440}\le n^{2+1/40-1/440}\le n^{2.0228}$; forming the lists costs $O(nd^2)=O(n^{1.05})$. So the $\tilde{O}(n/d)$ counting calls of Lemma~\ref{lem:realred}(a) cost
\[
\tilde{O}\big((n/d)\,n^{2.0229}\big)=\tilde{O}\big(n^{3-1/40+0.0229}\big)\leq O(n^{2.998})
\]
expected time, and its additional time is $\tilde{O}(n^3/d)=\tilde{O}(n^{2.975})$. 


\emph{$3$SUM.} By Corollary~\ref{cor:cmp}, a comparison count of Lemma~\ref{lem:realred}(b) costs
\[
O\Big(\frac{n^2}{d}\,n^{0.0229}+\frac{n^2}{d^{0.126}}+n^{2-1/440}\log n\Big)=O\big(n^{2-1/40+0.0229}\log n\big)=O(n^{1.9979}\log n),
\]
so $3$SUM can also be solved in $\tilde{O}(n^{1.9979})\leq O(n^{1.998})$ expected time.

\emph{Las Vegas, and high probability.} The reductions of~\cite{CVX22} never return a wrong answer and Corollary~\ref{cor:cmp} is deterministic, so the output is always correct, and only the running time is random. The bounds also hold with probability $1-n^{-c}$ for any constant $c$: we stop a run that exceeds twice the bound on its expected time and start again, and $c\log_2n$ runs all fail with probability at most $n^{-c}$.
\end{proof}

\subsection{Zero-Weight, Min-Weight, and Max-Weight \texorpdfstring{$k$}{k}-Clique}\label{sec:kclique}

The classical reduction of Ne\v{s}et\v{r}il and Poljak~\cite{NP85} from $k$-Clique to triangle detection turns the weighted $k$-Clique problems into Exact Triangle instances with $n^{\lfloor k/3\rfloor}$ vertices per part (after fixing $k-3\lfloor k/3\rfloor$ vertices when $3\nmid k$), so Theorem~\ref{thm:zwt} gives:
\begin{corollary}[Weighted $k$-Clique]\label{cor:kclique}
Let $k\ge3$ and $\nu\ge1$ be constants. Given a complete $k$-partite graph with parts of $n$ vertices and integer edge weights of absolute value at most $n^\nu$, deterministic algorithms decide in $O(n^{k-\varepsilon_T\lfloor k/3\rfloor})$ time whether some $k$-clique, with one vertex in each part, has total edge weight zero, and find a $k$-clique of minimum, or of maximum, total edge weight.
\end{corollary}
\begin{proof}
We apply the classical reduction from $k$-Clique to triangle detection of Ne\v{s}et\v{r}il and Poljak~\cite{NP85}. When $k$ is divisible by $3$, the reduction is simple. Split the $k$ parts into three groups of $k/3$ parts each, and build a tripartite graph $H$ whose vertices in the $i$th part (for $i=1,2,3$) are the $k/3$-cliques of the original graph with one vertex in each part of the $i$th group, with an edge between two such cliques when together they form a $2k/3$-clique. Since the original graph is complete $k$-partite, every choice of one vertex from each part of a group is a $k/3$-clique, and any two such cliques from different groups form a $2k/3$-clique, so $H$ is simply the complete tripartite graph with $n^{k/3}$ vertices per part. The triangles of $H$ correspond bijectively to the $k$-cliques of the original graph with one vertex in each part.

When $k$ is not divisible by $3$, let $t:=k-3\lfloor k/3\rfloor\in\{1,2\}$. We enumerate the $n^t$ ways of fixing one vertex in each of the first $t$ parts, and for each of them we build $H$ as above from the remaining $3\lfloor k/3\rfloor$ parts; now the triangles of $H$ correspond to the $k$-cliques through the $t$ fixed vertices. In both cases $H$ has $N:=n^{\lfloor k/3\rfloor}$ vertices per part, and we index its parts by $i\in\mathbb Z_3$.

We give $H$ edge weights so that a triangle and its $k$-clique have the same weight. Let $U$ be a vertex of $H$ in part $i$ and $V$ a vertex in part $i+1$. The weight of the edge $UV$ of $H$ is the total weight of the edges of the original graph between a vertex of $U$ and a vertex of $V$, of the edges inside $U$, of the edges between $U$ and the fixed vertices, and, if $i=1$ and $t=2$, of the edge between the two fixed vertices. The edges inside $V$ and the edges between $V$ and the fixed vertices are not counted here; they are counted in the edge from $V$ to part $i+2$. In this way every edge of a $k$-clique is counted exactly once over the three edges of its triangle, so there is no double-counting and the weights agree. The weight of an edge of $H$ has absolute value at most $\binom k2n^\nu\le N^{2\nu}$ for $n\ge k^2$.

Theorem~\ref{thm:zwt} decides whether $H$ has a triangle of weight zero in $O(N^{3-\varepsilon'}\log N)$ time, with the constants $\varepsilon'>\varepsilon_T$ of that theorem. By~\cite[Theorem~3.3]{VW13}, finding a maximum weight triangle costs $O(\log^2n)$ times as much, and a minimum weight triangle is a maximum weight triangle for the negated weights. Since $\varepsilon'>\varepsilon_T$, the logarithmic factors are absorbed, so each graph $H$ costs $O(N^{3-\varepsilon_T})$ time, and the $n^t$ graphs together cost $O(n^tN^{3-\varepsilon_T})=O(n^{k-\varepsilon_T\lfloor k/3\rfloor})$.
\end{proof}

Consequently, the Zero-Weight, Min-Weight, and Max-Weight $k$-Clique hypotheses~\cite{AL13,LVW18,BDT16} with polynomially bounded weights are refuted, also on general graphs, by the standard reduction to the complete $k$-partite form.

\subsection{Three conjectures of van den Brand, Nanongkai, and Saranurak}\label{sec:bns}

Van den Brand, Nanongkai, and Saranurak~\cite{BNS19} propose three ``hinted'' variants of the Online Matrix--Vector conjecture~\cite{HKNS15}, which they use to prove matching lower bounds for the running times of their dynamic matrix inverse algorithms. All three are over the Boolean semiring, allow polynomial preprocessing time in the first phase, and have parameters $t=n^{\tau}$ and $t_i=n^{\tau_i}$ for constants $\tau,\tau_i\in(0,1)$. We write $\omega(a,b,c)$ for the exponent of multiplying an $n^a\times n^b$ matrix by an $n^b\times n^c$ matrix. Each conjecture asserts that no algorithm {\em simultaneously} beats all of its listed bounds, for any $\varepsilon>0$.
\begin{itemize}
\item \emph{$v$-hinted Mv} (Definition~5.1 and Conjecture~5.2 of~\cite{BNS19}). Phase~1: an $n\times t$ matrix $M$. Phase~2: a $t\times n$ matrix $V$. Phase~3: an index $i$, after which the algorithm outputs $MV_{[n],i}$, the product of $M$ and column $i$ of $V$. Bounds conjectured to be impossible to achieve simultaneously: $n^{\omega(1,1,\tau)-\varepsilon}$ for Phase~2 and $n^{1+\tau-\varepsilon}$ for Phase~3.
\item \emph{Mv-hinted Mv} (Definition~5.6 and Conjecture~5.7  of~\cite{BNS19}). Phase~1: $N\in\{0,1\}^{n\times n}$ and $V\in\{0,1\}^{t\times n}$. Phase~2: a vector $I\in[n]^t$ of column indices. Phase~3: an index $j$, after which the algorithm outputs $N_{[n],I}V_{[t],j}$, where $N_{[n],I}$ is the $n\times t$ matrix whose $k$-th column is column $I_k$ of $N$. Bounds conjectured to be impossible to achieve simultaneously: $n^{\omega(1,\tau,1)-\varepsilon}$ for Phase~2 and $n^{1+\tau-\varepsilon}$ for Phase~3.
\item \emph{uMv-hinted uMv} (Definition~5.11 and Conjecture~5.12  of~\cite{BNS19}). Phase~1: $U\in\{0,1\}^{n\times t_1}$, $N\in\{0,1\}^{n\times n}$, and $V\in\{0,1\}^{t_2\times n}$. Phase~2: $I\in[n]^{t_1}$. Phase~3: $J\in[n]^{t_2}$. Phase~4: indices $i$ and $j$, after which the algorithm outputs $(UN_{I,J}V)_{i,j}$, where $N_{I,J}$ is the $t_1\times t_2$ submatrix of $N$ with the rows $I$ and the columns $J$. Bounds conjectured to be impossible to achieve simultaneously: $n^{\omega(1,\tau_1,1)-\varepsilon}$ for Phase~2, $n^{\omega(\tau_2,\tau_1,1)-\varepsilon}$ for Phase~3, and $n^{\tau_1+\tau_2-\varepsilon}$ for Phase~4.
\end{itemize}
The two obvious algorithms are to multiply everything out with fast matrix multiplication as soon as the hint arrives, or to wait and compute one matrix--vector product at the end, and the conjectures say that no algorithm is polynomially faster than both of these. The main lower bounds of~\cite{BNS19}, $n^{1.528}$ for dynamic matrix inverse with column updates and row queries and $n^{1.407}$ with element updates and element queries, use the conjectures at $\tau\approx0.53$ and at $\tau_1+\tau_2\approx1.41$, where they match the algorithms of~\cite{BNS19}. For thin hints, the conjectures are false, because the data structure of Corollary~\ref{cor:zt} is faster than both obvious algorithms.

\begin{corollary}[Hinted OMv]\label{cor:bns}
Conjectures~5.2 and~5.7 of~\cite{BNS19} fail for every $0<\tau<1/18$: Phase~2 takes $O(n^{2-0.063\tau})$ time and Phase~3 takes $O(n^{1+0.437\tau})$ time. Conjecture~5.12 fails for every $0<\tau_1<\tau_2/18$: Phase~3 takes $O(n^{1+\tau_2-0.063\tau_1})$ time and Phase~4 takes $O(n^{\tau_2+0.437\tau_1})$ time, with polynomial Phase~1 and a Phase~2 that only stores $I$.

More generally, Conjectures~5.2 and~5.7 fail for every $0<\tau<\varepsilon^*=0.1204\ldots$, and Conjecture~5.12 fails for every $0<\tau_1<\varepsilon^*\tau_2$. In each case the phases beat the conjectured bounds whatever the value of $\omega$, since $\omega(1,1,\tau)=\omega(1,\tau,1)\ge2$ and $\omega(\tau_2,\tau_1,1)\ge1+\tau_2$ because of the input and output sizes.
\end{corollary}
\begin{proof}
We compute over $\Z$ with $0/1$ matrices; a Boolean entry is 1 exactly when the corresponding integer entry is positive. Let $D:=t=n^\tau$ be the inner dimension. For $\tau<1/18$ we have $D^{18}=n^{18\tau}\le n$, so Corollary~\ref{cor:zt} applies with $N=n$.

\emph{Conjecture~5.2 of~\cite{BNS19}.} In Phase~2, we preprocess $X:=M$ and $Y:=V$ by Corollary~\ref{cor:zt}, in $O(n^2/D^{0.063})=O(n^{2-0.063\tau})$ time. In Phase~3, the $n$ entries of column $i$ of $MV$ are $n$ queries, which take $O(n\,D^{0.437})=O(n^{1+0.437\tau})$ time.

\emph{Conjecture~5.7 of~\cite{BNS19}.} The proof is the same, with $X:=N_{[n],I}$ and $Y:=V$, which are both known in Phase~2 ($X$ may have repeated columns, which do not affect the argument). Phase~1 only reads its input.

\emph{Conjecture~5.12 of~\cite{BNS19}.} Let $X:=U$, an $n\times t_1$ matrix, and $Y:=N_{I,J}$, a $t_1\times t_2$ matrix, which are both known in Phase~3, with the inner dimension $D:=t_1=n^{\tau_1}$. Cut the $n$ rows of $U$ into $n/t_2$ blocks of $t_2$ rows, and preprocess each of the $n/t_2$ products of a block with $Y$, which are products of a $t_2\times t_1$ matrix by a $t_1\times t_2$ matrix, by Corollary~\ref{cor:zt}. It applies because $t_2=n^{\tau_2}\ge n^{18\tau_1}=D^{18}$ when $\tau_1<\tau_2/18$. So Phase~3 costs $(n/t_2)\cdot O(t_2^2/D^{0.063})=O(n^{1+\tau_2-0.063\tau_1})$ time. In Phase~4, $(UN_{I,J}V)_{i,j}=\sum_{\ell=1}^{t_2}(UN_{I,J})_{i,\ell}V_{\ell,j}$ is a sum of at most $t_2$ entries of $XY$, which are $t_2$ queries and take $O(t_2D^{0.437})=O(n^{\tau_2+0.437\tau_1})$ time. Phase~2 only stores $I$.

\emph{General $\tau$.} Let $0<\tau<\varepsilon^*$, and fix $\varepsilon$ with $\tau<\varepsilon<\varepsilon^*$. By Corollary~\ref{cor:tradeoff}, with $c$ and $\theta$ chosen as in the proof of Theorem~\ref{thm:single} for $q:=\frac12$, there is, after absorbing the logarithmic factors, a $\gamma>0$ such that, since $D=n^\tau\le n^\varepsilon$, the arguments above go through with Phase~2 in $O(n^2/D^\gamma)=O(n^{2-\gamma\tau})$ time and Phase~3 in $O(n\,D^{1/2})=O(n^{1+\tau/2})$ time. For Conjecture~5.12, we use the same blocks, which requires $n^{\tau_1}\le n^{\varepsilon\tau_2}$, that is, $\tau_1\le\varepsilon\tau_2$ for some $\varepsilon<\varepsilon^*$.
\end{proof}

We conclude this section by discussing which results of~\cite{BNS19} rely on the conjectures in the refuted regime. From Conjectures~5.2 and~5.7,~\cite{BNS19} derive trade-offs between the worst-case update time $u$ and query time $q$ of dynamic algorithms, for every $0<\tau<1$: either $u=\Omega(n^{\omega(1,1,\tau)-\tau-\varepsilon})$ or $q=\Omega(n^{1+\tau-\varepsilon})$. These trade-offs cover the dynamic matrix product, inverse, and adjoint problems with column updates and row queries (Theorem~5.3 and Corollary~5.5 in~\cite{BNS19}) or with element updates and row queries (Theorem~5.8 and Corollary~5.9 in~\cite{BNS19}), and dynamic transitive closure and DAG path counting with vertex or edge updates and source queries (Corollaries~5.4 and~5.10 in~\cite{BNS19}). For element or pair queries, the condition becomes $q=\Omega(n^{\tau-\varepsilon})$ (Corollaries~5.9 and~5.10 in~\cite{BNS19}).

The main bounds, such as $u+q=\Omega(n^{1.528})$ for column updates and row queries, balance the two terms at $\tau\approx0.53$. For this setting, the current best upper bound on $\omega(1,1,\tau)$ is strictly greater than $2$, so the products are far from thin enough for our construction, and these bounds are still supported by the unrefuted forms of the conjectures. The same holds for the bounds from Conjecture~5.12 in~\cite{BNS19}, such as $\Omega(n^{1.407})$ for dynamic determinant, rank, and bipartite perfect matching (Corollaries~5.15 and~5.16 in~\cite{BNS19}), since Corollary~\ref{cor:bns} refutes Conjecture~5.12 in~\cite{BNS19} only for $\tau_1<0.1204\,\tau_2$, far from the $\tau_1+\tau_2\approx1.41$ at which these bounds are proved.

For $\tau<0.1204$, however, where $\omega(1,1,\tau)=2$, the trade-offs say that an update requires $n^{2-\tau-o(1)}$ time unless a row or source query takes $n^{1+\tau-o(1)}$ time, or, in the versions with element or pair queries, unless such a query takes $n^{\tau-o(1)}$ time; Corollary~\ref{cor:bns} refutes the forms of the conjectures that would imply these trade-offs. This is the end of the trade-offs with the fastest queries, where they match the algorithms of~\cite{San04,BNS19}, such as element updates in $O(n^{2-\tau})$ time with element queries in $O(n^{\tau})$ time~\cite[Theorem~3]{San04}. For small $\tau$, these algorithms are no longer known to be conditionally optimal.

The same goes for results that rely on this end of the trade-offs. Haeupler, Long, and Saranurak~\cite{HLS24} cite Corollary~5.10 of~\cite{BNS19} as showing that incremental and decremental distance oracles with approximation factor below $5/3$ cannot have worst-case update time $n^{2-\Omega(1)}$ and query time $n^{o(1)}$. With queries this fast, ruling out update time $n^{2-\delta}$ needs Conjecture~5.7 of~\cite{BNS19} for some $\tau<\delta$, so new evidence of hardness would be needed for $\delta\le0.1204$, although the remaining cases of the conjecture still rule out update time $O(n^{1.87})$. Van den Brand, Song, and Zhou~\cite{BSZ24} show that their data structure for dynamic attention, with amortized update time $n^{\omega(1,1,\tau)-\tau}$, is conditionally optimal for every $0<\tau\le1$ under a variant of Conjecture~5.7 of~\cite{BNS19} in which Phase~2 gives a matrix with at most $n^{\tau}$ nonzero entries instead of the vector $I$. Our proof applies to this variant as well, since after Phase~2 the product again has inner dimension at most $n^{\tau}$, so this lower bound needs a new hypothesis for $\tau<0.1204$.

